\scanpage{11}
\margin{livre de Lang\\ VA}
\textcircled{a}\quad Catégorie additive $\unsure{\mathscr A|\mathcal A}$

\textcircled{b}\quad Autodualité.
\[
 D:\mathscr A^{\unsure{o|\circ}}\longrightarrow\mathscr A
\]
(munie d'un isom
\[
 \unsure{k|\kappa}:\mathrm{id}_{\mathscr A}\longrightarrow D\circ D
\]
satisfaisant la condition de compatibilité habituelle).

\unsure{,} Soit alors
\[
 \unsure{\mathcal B|B}(A,B)=\mathrm{Hom}(A,D(B))
 \xrightarrow{\overset{s}{\sim}}\mathrm{Hom}(B,D(A))
\]
c'est un bifoncteur contravariant en $A,B$. On suppose de plus

\textcircled{c}\quad Pour tt $A\in\mathrm{Ob}\,\mathscr A$, un sous-groupe
\[
 \boxed{N(A)\subset B(A,A)^{\unsure{\text{antisymm}|\text{antisym}}}}
\]
dépendant de $A$ de façon\deleted{\illegible{1}} contravariante.
[N.B. $N$ n'est pas additif en $A$ ...]

Si $\xi\in N(A)$, on pose aussi $\varphi_\xi:A\longrightarrow D(A)$
le morphisme qu'il définit.

Soit $\xi\in N(A\times B)$, il définit
\[
 \varphi_{\xi}: A\times B\longrightarrow D(A)\times D(B)
\]
qui est antisym\uline{étrique}, donc \unsure{nécessairem|nécessairement} de
la forme
\[
 \begin{aligned}
 \begin{pmatrix}\alpha_\xi&{}^t\lambda_\xi\\-\lambda_\xi&\beta_\xi\end{pmatrix},
 &\quad \lambda_\xi\in\mathrm{Hom}(A,D(B))=B(A,B)\\
 &\quad \alpha_\xi\in B(A,A),\ \beta_\xi\in B(\deleted{\unsure{A}}B,B)
 \end{aligned}
\]
on obtient ainsi un hom
\[
 \boxed{\xi\leadsto\lambda_\xi:\ N(A\times B)\longrightarrow
 \mathrm{Hom}(A,D(B))=B(A,B)}
\]
% archive sheet number: 10

